%% chapter1.tex — Introduction
%%
%% Notes:
%%   - Use \chapter{Title} for the chapter heading.
%%   - Use \section, \subsection, \subsubsection for sub-headings.
%%   - Use \cite{key} for citations; they will appear in a "List of
%%     References" at the end of this chapter (chapterref mode) or in a
%%     single list after the last chapter (oneref mode).
%%   - Do NOT use \begin{document} or \end{document} here.

\chapter{Introduction}
\label{ch:introduction}

%% -----------------------------------------------------------------------
%% 1.1 Background
%% -----------------------------------------------------------------------
\section{Background and Motivation}
\label{sec:background}

Provide background information that motivates your research.
Explain the context of the problem and why it is important.
This section should be accessible to a reader who is knowledgeable in
the general field but may not be familiar with the specific topic.

A typical thesis introduction will survey the relevant literature~\cite{lamport1994latex}
and identify the gap in knowledge that this work addresses~\cite{knuth1986texbook}.
Each claim that draws on prior work should be supported by a citation.

%% -----------------------------------------------------------------------
%% 1.2 Problem Statement
%% -----------------------------------------------------------------------
\section{Problem Statement}
\label{sec:problem}

State the specific research problem clearly and concisely.
The problem statement should flow naturally from the motivation
established in Section~\ref{sec:background}.

\begin{itemize}
  \item What is not yet known or understood?
  \item Why does that gap matter?
  \item What is the scope of this work?
\end{itemize}

%% -----------------------------------------------------------------------
%% 1.3 Objectives
%% -----------------------------------------------------------------------
\section{Research Objectives}
\label{sec:objectives}

The objectives of this thesis are:
\begin{enumerate}
  \item To demonstrate the structure of a URI thesis template.
  \item To show how figures, tables, and equations are included.
  \item To illustrate the citation and bibliography workflow.
\end{enumerate}

%% -----------------------------------------------------------------------
%% 1.4 Contributions
%% -----------------------------------------------------------------------
\section{Contributions}
\label{sec:contributions}

The main contributions of this work are summarized below.

\begin{enumerate}
  \item A compilable, fully annotated thesis template conforming to URI
        Graduate School requirements.
  \item Worked examples of figures, tables, equations, and cross-references.
  \item Documentation of the BibTeX workflow for both per-chapter and
        single-list reference modes.
\end{enumerate}

%% -----------------------------------------------------------------------
%% 1.5 Thesis Organization
%% -----------------------------------------------------------------------
\section{Thesis Organization}
\label{sec:organization}

The remainder of this thesis is organized as follows.
Chapter~\ref{ch:background} provides a detailed review of related work.
Chapter~\ref{ch:methods} describes the methodology and experimental setup.
Conclusions and directions for future work are presented in
Chapter~\ref{ch:conclusion}.
Supplementary material appears in the appendix.
